Finite-element surrogates are usually trained on solutions with a displacement loss. For
linear elastostatics we show that, for the surrogates we trained, the norm of this loss, not
the use of solutions, limited stress accuracy. Minimising the assembled discrete potential
energy is exactly supervised regression in the energy norm, without solutions: the energy
gap of a prediction is the strain energy of its stress error. The Euclidean norm of the
displacement loss under-weights the stiff modes that carry the strains. In pre-registered
comparisons, changing only the norm of the supervised loss cut a three-dimensional
transformer's median energy-norm error \numNlThreeNormE{}-fold and its median von Mises
error \numNlThreeNormVm{}-fold; in two dimensions it more than halved the von Mises error,
also on held-out instances. In three dimensions, the transformer trained on energy alone
came within a factor of \numNlThreeLabelWithin{} of the energy-norm-supervised one in median
energy-norm and von Mises error and beat the displacement-supervised one on all five metrics
in every seed, while a displacement-supervised graph network had the networks' lowest
displacement error but a von Mises error \numThreeMgnOverFreeVm{} times the label-free
transformer's. Without solutions, energy differences rank predictions of a load case by
energy-norm error, and the sign of the energy flags those worse than the zero field. Given
the stiffness matrix, a \numRitzThreeKMax{}-unknown Ritz projection was more accurate than
the label-free transformer, averaged over the seeds, on every metric but peak stress, where its
median error was more than twice as large. No network transferred to a substantially finer
mesh.
